\documentclass[11pt,letterpaper]{article}

% --- Paquetes Esenciales ---
\usepackage[utf8]{inputenc}
\usepackage[T1]{fontenc}
\usepackage[spanish,es-tabla,es-nodecimaldot]{babel}
\usepackage{amsmath,amssymb,amsfonts,amsthm}
\usepackage{geometry}
\geometry{top=2.5cm, bottom=2.5cm, left=2.5cm, right=2.5cm}
\usepackage{booktabs}
\usepackage{array}
\usepackage{graphicx}
\usepackage{xcolor}
\usepackage{fancyhdr}
\usepackage{hyperref}
\usepackage{microtype}
\usepackage{cite}
\usepackage{caption}
\usepackage{subcaption}
\usepackage{siunitx} % Estandarización de unidades físicas (IEEE)

% --- Colores Institucionales y de Énfasis ---
\definecolor{tecblue}{RGB}{0, 43, 91}
\definecolor{tecred}{RGB}{179, 35, 23}
\definecolor{grayborder}{RGB}{200, 200, 200}
\definecolor{lightgray}{RGB}{245, 247, 250}

% --- Configuración de Hipervínculos ---
\hypersetup{
    colorlinks=true,
    linkcolor=tecblue,
    citecolor=tecred,
    urlcolor=tecblue,
    pdfauthor={Mario Cerdas, Sebastian Quiros Fernandez, Ian Fuentes Nunez},
    pdftitle={Memoria de Calculo y Analisis Teorico Completo - Sistema de Comunicacion por Medios Opticos},
    pdfsubject={EL-2113 Circuitos Electricos en Corriente Continua - ITCR}
}

% --- Encabezados y Pies de Página ---
\pagestyle{fancy}
\fancyhf{}
\fancyhead[L]{\small \textcolor{tecblue}{\textbf{ITCR | Escuela de Ingeniería Electrónica}}}
\fancyhead[R]{\small \textcolor{gray}{EL-2113: Memoria de Cálculo}}
\fancyfoot[C]{\thepage}
\renewcommand{\headrulewidth}{0.8pt}
\renewcommand{\footrulewidth}{0.4pt}

% --- Título y Datos Formales ---
\title{
    \vspace{-1.0cm}
    \Huge \textbf{\textcolor{tecblue}{Memoria de Cálculo y Justificación Teórica de Ingeniería}}\\[0.3cm]
    \Large \textbf{Sistema de Comunicación por Medios Ópticos de 4 Bits}\\[0.2cm]
    \large Proyecto de Investigación Dirigida --- Curso EL-2113: Circuitos Eléctricos en Corriente Continua
}

\author{
    \textbf{Mario Cerdas}\footnote{Carné: 2022437519, Correo: \texttt{m.cerdas@estudiantec.cr}} \quad
    \textbf{Sebastián Quirós Fernández}\footnote{Carné: 2024073526, Correo: \texttt{s.quiros.4@estudiantec.cr}} \quad
    \textbf{Ian Fuentes Núñez}\footnote{Carné: 2025100912, Correo: \texttt{I.fuentes.1@estudiantec.cr}}\\[0.3cm]
    \textit{Profesor: Ing. Fabián Alberto Picado García}\\[0.1cm]
    \textit{Escuela de Ingeniería Electrónica, Instituto Tecnológico de Costa Rica (ITCR)}\\
    \textit{Sede Central de Cartago, II Semestre 2026}
}

\date{\textbf{Última Revisión:} Noviembre de 2026}

\begin{document}

\maketitle
\thispagestyle{empty}

\begin{abstract}
\noindent El presente documento constituye la \textbf{Memoria de Cálculo Oficial y Análisis Teórico Integral} para el diseño, dimensionamiento, modelado y validación experimental del prototipo de \emph{Sistema de Comunicación por Medios Ópticos} del curso EL-2113. El sistema procesa secuencias digitales de cuatro bits ($N=4$) provenientes de señales periódicas y aperiódicas cuantizadas en MATLAB y transmitidas mediante un microcontrolador ATmega328P. La cadena de conversión y acondicionamiento implementa una red pasiva en escalera R--2R ($R=\SI{10}{\kilo\ohm}$, $2R=\SI{20}{\kilo\ohm}$) acoplada al componente activo OpAmp 1A de un circuito integrado LM358N, aislada mediante un seguidor de tensión (OpAmp 1B) que suprime una atenuación por efecto de carga del \SI{98.52}{\percent}. La señal analógica excita un generador óptico basado en un diodo emisor de luz (LED) rojo difuso ($\lambda \approx \SI{650}{\nano\meter}$), canalizado a través de una fibra óptica plástica (POF) de polimetilmetacrilato (PMMA) de \SI{45}{\centi\meter} hacia una fotorresistencia de sulfuro de cadmio (LDR GL5528). La variación fotoconductora modula un temporizador NE555 en configuración astable, generando tonos audibles entre \SI{176.7}{\hertz} y \SI{3443.0}{\hertz} sobre un buzzer piezoeléctrico pasivo desacoplado en corriente continua mediante un filtro pasa-altos ($f_c = \SI{15.9}{\hertz}$). Se detallan todas las deducciones por leyes de Kirchhoff (LCK y LVK), relaciones cuánticas de Planck-Einstein, física óptica de reflexión interna total, respuestas transitorias RC y consideraciones normativas de alimentación de laboratorio (Req. A.1 a Q.18).
\end{abstract}

\vspace{0.5cm}
\tableofcontents
\newpage

% =============================================================================
\section{Introducción y Arquitectura Modular del Sistema}
% =============================================================================

El objetivo central de la investigación dirigida consiste en diseñar, modelar en entornos virtuales (LTspice y TinkerCAD) y validar en un prototipo físico tangible un sistema optoelectrónico capaz de transmitir datos analógicos cuantizados mediante luz guiada y reproducir la información recibida en formato acústico. 

Conforme a la especificación oficial del curso (\emph{Requerimientos A.1 a Q.18}), el sistema se estructura de forma estrictamente modular en siete etapas funcionales delimitadas e identificables (Req. A.1 y B.2), como se sintetiza en el Cuadro~\ref{tab:etapas_sistema}.

\begin{table}[htbp]
\centering
\caption{Delimitación de las Siete Etapas Funcionales del Sistema.}
\label{tab:etapas_sistema}
\small
\renewcommand{\arraystretch}{1.2}
\begin{tabular}{cp{4.8cm}p{4.2cm}p{4.5cm}}
\toprule
\textbf{Etapa} & \textbf{Denominación Funcional} & \textbf{Componentes Clave} & \textbf{Requerimientos Normativos} \\
\midrule
1 & Patrón de Datos de Prueba & Arduino UNO (ATmega328P) & C.3, D.4 (4 bits, \SI{1}{\milli\second} a \SI{10}{\second}, MATLAB) \\
2 & Convertidor D/A (DAC Pasivo + Activo) & Red R--2R (\SI{10}{\kilo\ohm} / \SI{20}{\kilo\ohm}) + OpAmp 1A & E.5, O.16 (Pasivos + 1 Activo, Error $< \SI{10}{\percent}$) \\
3 & Interfaz de Aislamiento & Seguidor OpAmp 1B (LM358N) & F.6, B.2 (Solo OpAmp, supresión de carga) \\
4 & Generador Óptico & Resistor \SI{150}{\ohm} + LED Rojo (\SI{650}{\nano\meter}) & G.7, H.8 (Solo pasivos + luz no láser) \\
5 & Medio y Receptor Óptico Pasivo & Fibra POF PMMA (\SI{45}{\centi\meter}) + LDR CdS & I.9, J.10 (30--\SI{60}{\centi\meter}, receptor 100\% pasivo) \\
6 & Generador de Onda Cuadrada & Temporizador NE555 astable & K.11 (CI activo comercial + pasivos) \\
7 & Transductor Acústico Final & Condensador \SI{10}{\micro\farad} + Buzzer Pasivo & L.12 (Solo buzzer piezoeléctrico pasivo) \\
\bottomrule
\end{tabular}
\end{table}

% =============================================================================
\section{Etapa 1: Patrón de Datos de Prueba (MATLAB y Arduino)}
% =============================================================================

\subsection{Fundamento Teórico de Cuantización Uniforme de 4 Bits}
De acuerdo con la teoría clásica de procesamiento digital de señales \cite{proakis2007,sedra2020}, la discretización en amplitud de una señal continua $x(t)$ acotada en el intervalo dinámico $[0, V_{\text{ref}}]$ con $V_{\text{ref}} = \SI{5.0000}{\volt}$ mediante una palabra binaria de $N = 4$ bits genera un espacio finito de estados cuantizados:
\begin{equation}
M = 2^N = 2^4 = 16\text{ niveles discretos}\quad (D \in \{0, 1, 2, \dots, 15\})
\label{eq:niveles}
\end{equation}

El incremento uniforme de tensión correspondiente al bit menos significativo ($\Delta V_{\text{LSB}}$) representa la resolución analógica mínima del sistema:
\begin{equation}
\Delta V_{\text{LSB}} = \frac{V_{\text{ref}}}{2^N} = \frac{\SI{5.0000}{\volt}}{16} = \SI{0.3125}{\volt} = \SI{312.5}{\milli\volt}
\label{eq:lsb_teorico}
\end{equation}

La asignación del código entero de 4 bits para una muestra temporal arbitraria $x(t_k)$ se define formalmente mediante el operador de redondeo simétrico y saturación de rango:
\begin{equation}
D[k] = \mathrm{sat}_{[0,15]} \left( \mathrm{round}\left( \frac{x(t_k) - V_{\text{min}}}{\Delta V_{\text{LSB}}} \right) \right), \quad k = 0, 1, \dots, L-1
\label{eq:cuantizador}
\end{equation}
donde cada código $D[k]$ se proyecta en su descomposición posicional binaria natural $(b_3 b_2 b_1 b_0)_2$:
\begin{equation}
D[k] = \sum_{i=0}^{3} b_i [k] \cdot 2^i = 8\,b_3[k] + 4\,b_2[k] + 2\,b_1[k] + 1\,b_0[k], \quad b_i \in \{0, 1\}
\label{eq:binario_posicional}
\end{equation}

\subsection{Relación Señal a Ruido de Cuantización Teórica (SQNR)}
Para una señal armónica continua senoidal que recorre la totalidad de la escala del conversor, la potencia teórica de la señal es $P_s = (V_{\text{FS}} / (2\sqrt{2}))^2 = V_{\text{FS}}^2 / 8$, y la varianza del error de cuantización bajo la hipótesis de ruido blanco uniforme distribuido en $[-\Delta V/2, \Delta V/2]$ es $\sigma_e^2 = \Delta V^2 / 12$. La relación señal a ruido de cuantización teórica (SQNR) se determina como \cite{proakis2007}:
\begin{equation}
\mathrm{SQNR} = 10\log_{10}\left( \frac{P_s}{\sigma_e^2} \right) = 6.02\,N + \SI{1.76}{\decibel}
\label{eq:sqnr_formula}
\end{equation}
Para $N = 4$ bits:
\begin{equation}
\mathrm{SQNR} = 6.02(4) + 1.76 = 24.08 + 1.76 = \SI{25.84}{\decibel}
\end{equation}

\subsection{Síntesis de Señales en MATLAB y Arreglos C++ (Req. D.4)}
En el script \texttt{Matlab/Primero.m} se implementan dos funciones de excitación para cumplir el requerimiento \textbf{D.4}:
\begin{enumerate}
    \item \textbf{Señal Periódica ($L = 20$ muestras):} Onda senoidal pura definida por $x_p(t) = 2.34 + 2.10\sin(2\pi f_0 t)\,\si{\volt}$, acotada entre \SI{0.24}{\volt} y \SI{4.44}{\volt}, generando el arreglo:
    \begin{equation*}
    \texttt{PATRON\_PERIODICO} = \{7, 10, 12, 13, 14, 14, 14, 12, 11, 9, 6, 4, 3, 1, 1, 1, 2, 3, 5, 7\}
    \end{equation*}
    \item \textbf{Señal Aperiódica Transitoria ($L = 18$ muestras):} Función amortiguada exponencialmente representativa de un transitorio real $x_a(t) = 0.40 + 4.25 e^{-3.2 t} |\sin(3\pi t + 0.3)|\,\si{\volt}$, produciendo el arreglo:
    \begin{equation*}
    \texttt{PATRON\_APERIODICO} = \{5, 10, 10, 8, 5, 2, 3, 4, 4, 3, 2, 1, 2, 2, 2, 2, 1, 1\}
    \end{equation*}
\end{enumerate}

\subsection{Temporización No Bloqueante en Firmware Arduino (Req. C.3)}
El firmware \texttt{Arduino/Logica/Logica.ino} transfiere secuencialmente los bits a los terminales digitales de salida del ATmega328P:
\begin{equation*}
\texttt{D8} \to b_0\text{ (LSB)}, \quad \texttt{D9} \to b_1, \quad \texttt{D10} \to b_2, \quad \texttt{D11} \to b_3\text{ (MSB)}
\end{equation*}
Para satisfacer el requerimiento \textbf{C.3}, cada palabra binaria se sostiene durante un periodo de retención programable $T_{\text{est}}$ gobernado por el temporizador por hardware del microcontrolador mediante llamadas a \texttt{millis()}, asegurando:
\begin{equation}
\SI{1}{\milli\second} \le T_{\text{est}} \le \SI{10}{\second} \quad (\SI{1}{\milli\second} \le T_{\text{est}} \le \SI{10000}{\milli\second})
\label{eq:test_req}
\end{equation}
Esta arquitectura previene el bloqueo de la CPU inducido por funciones primitivas como \texttt{delay()}, posibilitando la lectura asíncrona de comandos por el puerto serie UART a 9600 baudios.

% =============================================================================
\section{Etapa 2: Convertidor Digital-Analógico (DAC R--2R y OpAmp 1A)}
% =============================================================================

\subsection{Justificación Técnica de la Topología en Escalera R--2R (Req. E.5)}
El requerimiento \textbf{E.5} exige implementar el conversor con elementos pasivos y un elemento activo. Se evaluaron las dos arquitecturas clásicas de conversión D/A \cite{franco2015,sedra2020}:

\begin{table}[htbp]
\centering
\caption{Comparación Técnica de Topologías de Conversión D/A.}
\label{tab:comparacion_dac}
\small
\renewcommand{\arraystretch}{1.2}
\begin{tabular}{p{4.5cm}p{5.5cm}p{5.5cm}}
\toprule
\textbf{Parámetro de Diseño} & \textbf{DAC Resistencias Ponderadas ($R, 2R, 4R, 8R$)} & \textbf{DAC Escalera R--2R (Seleccionado)} \\
\midrule
Dispersión de componentes & Requiere 4 valores distintos ($10, 20, 40, 80\text{ k}\Omega$). Difícil emparejamiento. & Requiere únicamente \textbf{2 valores estándar}: $R=\SI{10}{\kilo\ohm}$ y $2R=\SI{20}{\kilo\ohm}$ (precisión al 1\%). \\
Impedancia equivalente ($R_{\text{th}}$) & Variable en función del código digital aplicado. & \textbf{Rigurosamente constante}: $R_{\text{th}} = R = \SI{10}{\kilo\ohm}$ para los 16 códigos. \\
Alimentación del OpAmp & Sumador inversor genera tensión negativa ($V_{\text{out}} < \SI{0}{\volt}$). Exige fuente dual $\pm V_{\text{CC}}$. & Seguidor unitario opera con tensión positiva en fuente unipolar referenciada a tierra (\SI{0}{\volt}). \\
Sensibilidad a tolerancias & Alta: el bit MSB domina fuertemente y provoca saltos no monótonos. & Baja: la relación $2:1$ garantiza monotonicidad estricta y error $< \SI{0.002}{\percent} \ll \SI{10}{\percent}$. \\
\bottomrule
\end{tabular}
\end{table}

Por consiguiente, se selecciona la \textbf{escalera R--2R} acoplada al componente activo \textbf{OpAmp 1A} (Pines 1, 2 y 3 del circuito integrado LM358N).

\subsection{Deducción Matemática Rigurosa por Análisis Nodal LCK}
Considérese el diagrama esquemático de la escalera R--2R de 4 bits con terminación a tierra:

\begin{figure}[htbp]
\centering
\small
\begin{verbatim}
       V_b0 (D8)        V_b1 (D9)        V_b2 (D10)       V_b3 (D11)
         │                │                │                │
       [ 2R ]           [ 2R ]           [ 2R ]           [ 2R ]
         │                │                │                │
     V1  ●─────[ R ]──────●─────[ R ]──────●─────[ R ]──────● V4 = V_DAC
         │               V2               V3                │
       [ 2R ] (GND)                                         └──► (+) OpAmp 1A
         │
        GND
\end{verbatim}
\caption{Diagrama circuital de la escalera R--2R de 4 bits.}
\label{fig:diagrama_r2r}
\end{figure}

Sean $V_{b0}, V_{b1}, V_{b2}, V_{b3} \in \{\SI{0}{\volt}, \SI{5}{\volt}\}$ los niveles de tensión aplicados por el microcontrolador ($V_{bi} = V_{\text{ref}} \cdot b_i$). Asumiendo que la entrada no inversora del OpAmp presenta impedancia idealmente infinita ($I_{\text{in+}} = \SI{0}{\ampere}$), se aplica la \textbf{Ley de Corrientes de Kirchhoff (LCK)} en cada nodo ($\sum I_{\text{salientes}} = 0$) \cite{hayt2012,alexander2018}:

\paragraph{1. Ecuación Nodal en el Nodo 1 ($V_1$):}
\begin{equation}
\frac{V_1 - 0}{2R} + \frac{V_1 - V_{b0}}{2R} + \frac{V_1 - V_2}{R} = 0
\label{eq:lck_nodo1}
\end{equation}
Multiplicando la ecuación (\ref{eq:lck_nodo1}) por $2R$:
\begin{equation}
V_1 + (V_1 - V_{b0}) + 2(V_1 - V_2) = 0 \implies 4V_1 - 2V_2 = V_{b0}
\end{equation}
Despejando $V_1$:
\begin{equation}
V_1 = \frac{V_2}{2} + \frac{V_{b0}}{4} \label{eq:v1_despejado}
\end{equation}

\paragraph{2. Ecuación Nodal en el Nodo 2 ($V_2$):}
\begin{equation}
\frac{V_2 - V_1}{R} + \frac{V_2 - V_{b1}}{2R} + \frac{V_2 - V_3}{R} = 0
\label{eq:lck_nodo2}
\end{equation}
Multiplicando por $2R$:
\begin{equation}
2(V_2 - V_1) + (V_2 - V_{b1}) + 2(V_2 - V_3) = 0 \implies -2V_1 + 5V_2 - 2V_3 = V_{b1}
\end{equation}
Sustituyendo la ecuación (\ref{eq:v1_despejado}) en la expresión anterior:
\begin{equation}
-2\left( \frac{V_2}{2} + \frac{V_{b0}}{4} \right) + 5V_2 - 2V_3 = V_{b1} \implies -V_2 - \frac{V_{b0}}{2} + 5V_2 - 2V_3 = V_{b1}
\end{equation}
\begin{equation}
4V_2 - 2V_3 = V_{b1} + \frac{V_{b0}}{2} \implies V_2 = \frac{V_3}{2} + \frac{V_{b1}}{4} + \frac{V_{b0}}{8} \label{eq:v2_despejado}
\end{equation}

\paragraph{3. Ecuación Nodal en el Nodo 3 ($V_3$):}
\begin{equation}
\frac{V_3 - V_2}{R} + \frac{V_3 - V_{b2}}{2R} + \frac{V_3 - V_4}{R} = 0
\label{eq:lck_nodo3}
\end{equation}
Multiplicando por $2R$:
\begin{equation}
2(V_3 - V_2) + (V_3 - V_{b2}) + 2(V_3 - V_4) = 0 \implies -2V_2 + 5V_3 - 2V_4 = V_{b2}
\end{equation}
Sustituyendo la ecuación (\ref{eq:v2_despejado}):
\begin{equation}
-2\left( \frac{V_3}{2} + \frac{V_{b1}}{4} + \frac{V_{b0}}{8} \right) + 5V_3 - 2V_4 = V_{b2} \implies -V_3 - \frac{V_{b1}}{2} - \frac{V_{b0}}{4} + 5V_3 - 2V_4 = V_{b2}
\end{equation}
\begin{equation}
4V_3 - 2V_4 = V_{b2} + \frac{V_{b1}}{2} + \frac{V_{b0}}{4} \implies V_3 = \frac{V_4}{2} + \frac{V_{b2}}{4} + \frac{V_{b1}}{8} + \frac{V_{b0}}{16} \label{eq:v3_despejado}
\end{equation}

\paragraph{4. Ecuación Nodal en el Nodo de Salida ($V_4 = V_{\text{DAC}}$):}
En el nodo 4, dado que la corriente hacia la entrada no inversora del OpAmp es idealmente nula ($I_{\text{in+}} = 0$):
\begin{equation}
\frac{V_4 - V_3}{R} + \frac{V_4 - V_{b3}}{2R} = 0
\label{eq:lck_nodo4}
\end{equation}
Multiplicando por $2R$:
\begin{equation}
2(V_4 - V_3) + (V_4 - V_{b3}) = 0 \implies -2V_3 + 3V_4 = V_{b3}
\end{equation}
Sustituyendo la ecuación (\ref{eq:v3_despejado}):
\begin{equation}
-2\left( \frac{V_4}{2} + \frac{V_{b2}}{4} + \frac{V_{b1}}{8} + \frac{V_{b0}}{16} \right) + 3V_4 = V_{b3}
\end{equation}
\begin{equation}
-V_4 - \frac{V_{b2}}{2} - \frac{V_{b1}}{4} - \frac{V_{b0}}{8} + 3V_4 = V_{b3} \implies 2V_4 = V_{b3} + \frac{V_{b2}}{2} + \frac{V_{b1}}{4} + \frac{V_{b0}}{8}
\end{equation}
Dividiendo la igualdad entre 2:
\begin{equation}
V_4 = \frac{V_{b3}}{2} + \frac{V_{b2}}{4} + \frac{V_{b1}}{8} + \frac{V_{b0}}{16}
\label{eq:v4_fracciones}
\end{equation}
Sustituyendo $V_{bi} = V_{\text{ref}} \cdot b_i$:
\begin{equation}
\boxed{V_{\text{DAC}} = V_4 = V_{\text{ref}} \left( \frac{b_3}{2} + \frac{b_2}{4} + \frac{b_1}{8} + \frac{b_0}{16} \right) = \frac{V_{\text{ref}}}{16} \sum_{i=0}^{3} b_i 2^i = V_{\text{ref}} \cdot \frac{D}{16}}
\label{eq:vdac_final}
\end{equation}

\subsection{Demostración de la Resistencia Equivalente de Thévenin ($R_{\text{th}}$)}
Se demuestra por reducción iterativa de fuentes pasivadas (cortocircuito a tierra) \cite{hayt2012}:
\begin{enumerate}
    \item En el nodo 1: La resistencia de terminación $2R$ en paralelo con la rama vertical $2R$ produce:
    \begin{equation}
    R_{\text{eq1}} = 2R \parallel 2R = \frac{2R \cdot 2R}{2R + 2R} = R
    \end{equation}
    \item En el nodo 2: $R_{\text{eq1}}$ en serie con la resistencia horizontal $R$ resulta en $2R$. Al combinarse en paralelo con la rama vertical del bit 1 ($2R$):
    \begin{equation}
    R_{\text{eq2}} = (R + R) \parallel 2R = 2R \parallel 2R = R
    \end{equation}
    \item En el nodo 3: Repitiendo el proceso inductivo:
    \begin{equation}
    R_{\text{eq3}} = (R + R) \parallel 2R = 2R \parallel 2R = R
    \end{equation}
    \item En el nodo 4: Finalmente:
    \begin{equation}
    \boxed{R_{\text{th}} = (R + R) \parallel 2R = 2R \parallel 2R = R = \SI{10.0}{\kilo\ohm}}
    \end{equation}
\end{enumerate}
La resistencia equivalente es \textbf{estrictamente invariante frente a cualquier combinación de código digital}.

\subsection{Criterio de Selección de $R=\SI{10}{\kilo\ohm}$ y Corriente de Pines}
La resistencia vista por cualquier pin digital conectado a una rama $2R$ satisface $R_{\text{in}} \ge 2R$ \cite{franco2015}. Por Ley de Ohm:
\begin{equation}
I_{\text{pin(max)}} = \frac{V_{\text{ref}}}{2R} = \frac{\SI{5.0000}{\volt}}{\SI{20.0}{\kilo\ohm}} = \SI{0.250}{\milli\ampere} = \SI{250}{\micro\ampere} \ll \SI{20}{\milli\ampere}
\label{eq:corriente_pin}
\end{equation}
La resistencia de conducción interna de los MOSFETs de salida del puerto digital del ATmega328P es $R_{\text{on}} \approx \SI{25}{\ohm}$ \cite{sedra2020}. La caída interna inducida es apenas:
\begin{equation}
\Delta V_{\text{drop}} = I_{\text{pin(max)}} \cdot R_{\text{on}} = \SI{0.250}{\milli\ampere} \times \SI{25}{\ohm} = \SI{6.25}{\milli\volt}
\end{equation}
lo que introduce un error relativo despreciable de:
\begin{equation}
E_{\text{CMOS}} = \frac{\SI{6.25}{\milli\volt}}{\SI{5000}{\milli\volt}} \times \SI{100}{\percent} = \SI{0.125}{\percent} \ll \SI{10}{\percent}
\end{equation}
garantizando la integridad del nivel lógico $V_{\text{OH}} \approx \SI{5.00}{\volt}$.

\subsection{Polarización del LM358N y Margen Dinámico Superior}
El amplificador operacional comercial LM358N posee una etapa de salida cuasi-complementaria con saturación alta limitada por \cite{ti_lm358}:
\begin{equation}
V_{\text{OH(max)}} \le V_{\text{CC}} - \SI{1.50}{\volt}
\label{eq:voh_lm358}
\end{equation}
Si el circuito integrado se polarizara con la fuente de \SI{+5.00}{\volt}, el riel superior alcanzaría únicamente $V_{\text{OH}} \approx \SI{3.50}{\volt}$, provocando una distorsión por recorte severo en los códigos $D = 12, 13, 14, 15$ (\SI{3.75}{\volt} a \SI{4.6875}{\volt}). Para asegurar excursión analógica totalmente lineal, el terminal $V_{\text{CC}}$ (Pin 8) se alimenta con \SI{+9.00}{\volt} desde el segundo canal regulado de la fuente de laboratorio:
\begin{equation}
V_{\text{OH(disponible)}} = \SI{9.00}{\volt} - \SI{1.50}{\volt} = \SI{7.50}{\volt} > \SI{4.6875}{\volt} \quad (\text{Margen de seguridad: } \SI{+2.81}{\volt})
\end{equation}

% =============================================================================
\section{Etapa 3: Interfaz de Aislamiento y Efecto de Carga}
% =============================================================================

\subsection{Demostración Analítica del Efecto de Carga (Loading Effect)}
Para evidenciar la necesidad obligatoria de la Etapa 3 conforme a los requerimientos \textbf{B.2} y \textbf{F.6}, considérese el escenario hipotético de conexión directa de la red R--2R a la malla del diodo emisor ($R_{\text{driver}} = \SI{150}{\ohm}$) sin acoplamiento activo:

\begin{figure}[htbp]
\centering
\small
\begin{verbatim}
     V_DAC o──────[ R_th = 10 kΩ ]──────o V_carga
                                        │
                                      [ R_driver = 150 Ω ]
                                        │
                                       (>|) D_LED (V_D ≈ 1.8 V)
                                        │
                                       GND
\end{verbatim}
\caption{Circuito equivalente de Thévenin con acoplamiento directo de carga.}
\label{fig:loading_effect}
\end{figure}

Por el teorema de divisor de tensión resistivo \cite{alexander2018}:
\begin{equation}
V_{\text{carga}} = V_{\text{DAC}} \cdot \frac{R_{\text{driver}}}{R_{\text{th}} + R_{\text{driver}}} = V_{\text{DAC}} \cdot \frac{\SI{150}{\ohm}}{\SI{10000}{\ohm} + \SI{150}{\ohm}} = \frac{150}{10\,150}\,V_{\text{DAC}} \approx 0.014778\,V_{\text{DAC}}
\label{eq:divisor_carga}
\end{equation}
La atenuación o porcentaje de pérdida de amplitud analógica es:
\begin{equation}
\text{Pérdida} = \left( 1 - 0.014778 \right) \times \SI{100}{\percent} = \mathbf{\SI{98.52}{\percent}}
\label{eq:perdida_carga}
\end{equation}
Bajo esta condición, en el estado máximo ($D=15$, $V_{\text{DAC}} = \SI{4.6875}{\volt}$), la tensión resultante sobre la carga sería de únicamente:
\begin{equation}
V_{\text{carga(max)}} = \SI{4.6875}{\volt} \times 0.014778 = \SI{0.0693}{\volt} = \SI{69.3}{\milli\volt}
\end{equation}
Como $\SI{69.3}{\milli\volt} \ll V_D \approx \SI{1.80}{\volt}$, el diodo emisor \textbf{permanecería perpetuamente apagado}, haciendo imposible la comunicación óptica.

\subsection{Aislamiento Unipolar mediante Seguidor de Tensión (OpAmp 1B)}
Al intercalar el \textbf{OpAmp 1B} (Pines 5, 6 y 7 del LM358N) en configuración buffer seguidor ($A_v = 1$, con lazo de realimentación cerrado entre Pin 6 y Pin 7) \cite{franco2015}:
\begin{itemize}
    \item \textbf{Impedancia de entrada:} $Z_{\text{in}} > \SI{1e11}{\ohm} \implies I_{\text{entrada}} \to \SI{0}{\ampere}$ (no extrae corriente ni atenúa la escalera R--2R).
    \item \textbf{Impedancia de salida:} $Z_{\text{out}} = \frac{R_o}{1 + A_{\text{OL}}\beta} \approx \SI{0}{\ohm}$ (suministra la corriente requerida por el diodo sin caída de tensión).
    \item \textbf{Tensión transferida:} $V_{\text{interfaz}} = V_{\text{DAC}} = \frac{V_{\text{ref}}}{16} \cdot D$.
\end{itemize}

% =============================================================================
\section{Etapa 4: Generador Óptico y Justificación Física del Color}
% =============================================================================

\subsection{Cálculo del Resistor Limitador de Corriente ($R_{\text{driver}}$ — Req. G.7)}
Aplicando la Ley de Voltajes de Kirchhoff (LVK) en la malla del emisor:
\begin{equation}
-V_{\text{interfaz}}(D) + I_{\text{LED}}(D) \cdot R_{\text{driver}} + V_D = 0
\end{equation}
Despejando la corriente de excitación:
\begin{equation}
I_{\text{LED}}(D) = \begin{cases} 
\SI{0}{\ampere}, & V_{\text{interfaz}}(D) \le V_D \\ 
\dfrac{V_{\text{interfaz}}(D) - V_D}{R_{\text{driver}}}, & V_{\text{interfaz}}(D) > V_D 
\end{cases}
\label{eq:lvk_corriente_led}
\end{equation}

Para fijar una corriente de operación máxima de seguridad $I_{\text{LED(max)}} \approx \SI{19.25}{\milli\ampere} < \SI{20.0}{\milli\ampere}$ en el código supremo ($D = 15$, $V_{\text{interfaz}} = \SI{4.6875}{\volt}$) considerando $V_D = \SI{1.8000}{\volt}$ para un LED rojo difuso:
\begin{equation}
R_{\text{driver}} = \frac{V_{\text{interfaz(15)}} - V_D}{I_{\text{LED(max)}}} = \frac{\SI{4.6875}{\volt} - \SI{1.8000}{\volt}}{\SI{0.01925}{\ampere}} = \frac{\SI{2.8875}{\volt}}{\SI{0.01925}{\ampere}} = \mathbf{\SI{150.0}{\ohm}}
\label{eq:rdriver_calculo}
\end{equation}
El valor resultante coincide exactamente con un valor comercial estándar al \SI{5}{\percent} de la serie E12/E24 (\SI{150}{\ohm}). La potencia máxima disipada por el componente pasivo es:
\begin{equation}
P_{R_{\text{driver}}} = I_{\text{LED(max)}}^2 \cdot R_{\text{driver}} = (\SI{0.01925}{\ampere})^2 \times \SI{150}{\ohm} = \SI{0.05558}{\watt} = \mathbf{\SI{55.6}{\milli\watt}} \ll \SI{250}{\milli\watt}\text{ (1/4 W)}
\label{eq:potencia_rdriver}
\end{equation}
El resistor trabaja al \SI{22.2}{\percent} de su capacidad nominal, previniendo derivas térmicas por sobrecalentamiento.

\subsection{Justificación Teórica del Color Rojo ($\lambda \approx \SI{650}{\nano\meter}$) basada en Física 2 y 3 (Req. H.8)}
El requerimiento \textbf{H.8} establece plena libertad en la elección de la longitud de onda de emisión, exigiendo justificación fundamentada en física. La elección del espectro rojo visible ($\lambda \approx \SI{650}{\nano\meter}$) sobre otras opciones cromáticas (verde, azul o infrarrojo) se apoya en tres principios físicos:

\paragraph{1. Relación Cuántica de Planck-Einstein y Margen Dinámico Activo:}
Por la física cuántica del estado sólido \cite{tipler2010,kasap2013}, la recombinación radiativa electrón-hueco en la juntura semiconductora libera fotones cuya energía fundamental obedece:
\begin{equation}
E = h\nu = \frac{h \cdot c}{\lambda} \approx q \cdot V_D \implies V_D \approx \frac{h \cdot c}{q \cdot \lambda}
\label{eq:planck_einstein}
\end{equation}
donde las constantes físicas fundamentales son:
\begin{itemize}
    \item Constante de Planck: $h = \SI{6.62607015e-34}{\joule\second}$
    \item Velocidad de la luz en el vacío: $c = \SI{2.99792458e8}{\meter\per\second}$
    \item Carga elemental: $q = \SI{1.602176634e-19}{\coulomb}$
\end{itemize}

Evaluando las distintas regiones del espectro visible:
\begin{itemize}
    \item \textbf{LED Rojo ($\lambda \approx \SI{650}{\nano\meter}$):}
    \begin{equation}
    E = \frac{(\num{6.626e-34})(\num{3.00e8})}{(\num{1.602e-19})(\num{650e-9})} = \SI{1.908}{\electronvolt} \implies V_D \approx \SI{1.80}{\volt}
    \end{equation}
    Conduce a partir del código $D=6$ ($V_6 = \SI{1.875}{\volt}$), garantizando \textbf{10 estados activos continuos de modulación óptica} ($D=6$ a $D=15$).
    \item \textbf{LED Verde ($\lambda \approx \SI{525}{\nano\meter}$):} $E = \SI{2.36}{\electronvolt} \implies V_D \approx \SI{2.40}{\volt}$. Comienza a conducir a partir de $D=8$ (\SI{2.50}{\volt}), desaprovechando el \SI{50}{\percent} de los códigos digitales.
    \item \textbf{LED Azul ($\lambda \approx \SI{460}{\nano\meter}$):} $E = \SI{2.70}{\electronvolt} \implies V_D \approx \SI{3.0}{\volt}\text{--}\SI{3.2}{\volt}$. Conduce únicamente desde $D=11$ (\SI{3.4375}{\volt}), recortando el \SI{68.7}{\percent} del rango dinámico.
\end{itemize}

\paragraph{2. Ventana de Mínima Atenuación en Fibras Plásticas de PMMA:}
La fibra óptica plástica (POF) emplea un núcleo de polimetilmetacrilato ($n_1 = 1.492$) y revestimiento fluorado ($n_2 = 1.402$). Por la Ley de Snell \cite{tipler2010,ziemann2008}, el ángulo crítico de Reflexión Interna Total (TIR) y la Apertura Numérica (NA) son:
\begin{equation}
\theta_c = \arcsin\left( \frac{n_2}{n_1} \right) = \arcsin\left( \frac{1.402}{1.492} \right) = \SI{69.99}{\degree} \approx \SI{70.0}{\degree}
\label{eq:angulo_critico}
\end{equation}
\begin{equation}
\mathrm{NA} = \sqrt{n_1^2 - n_2^2} = \sqrt{(1.492)^2 - (1.402)^2} = \sqrt{2.2261 - 1.9656} = \sqrt{0.2605} \approx \mathbf{0.510}
\label{eq:apertura_numerica}
\end{equation}

El espectro de pérdidas ópticas en polímeros PMMA está condicionado por dos fenómenos físicos contrapuestos \cite{ziemann2008}:
\begin{enumerate}
    \item \textbf{Dispersión intrínseca de Rayleigh ($\alpha_{\text{Rayleigh}} \propto \lambda^{-4}$):} Domina fuertemente en longitudes de onda cortas (azul y ultravioleta).
    \item \textbf{Absorción molecular por sobretonos vibracionales C--H:} Crece exponencialmente a partir de $\lambda > \SI{700}{\nano\meter}$ (infrarrojo cercano).
\end{enumerate}
La superposición de ambos mecanismos define la célebre \textbf{ventana de transmisión de mínima atenuación del PMMA centrada exactamente en $\lambda = \SI{650}{\nano\meter}$}, con coeficientes de pérdida de apenas $\alpha \approx \SI{0.15}{\decibel\per\meter}$ \cite{ziemann2008}.

\paragraph{3. Acoplamiento Espectral con el Sulfuro de Cadmio (CdS):}
La fotorresistencia está constituida por CdS policristalino ($E_g \approx \SI{2.42}{\electronvolt}$). Sus niveles de defecto concentran la máxima fotosensibilidad en la banda de \SI{540}{\nano\meter} a \SI{650}{\nano\meter} \cite{rose1963,streetman2016}. Mientras que emisores infrarrojos (\SI{850}{\nano\meter}--\SI{940}{\nano\meter}, $E \approx \SI{1.3}{\electronvolt} < E_g$) no poseen suficiente energía fotónica para promover portadores, el LED rojo visible genera la máxima variación de conductancia.

% =============================================================================
\section{Etapa 5: Medio Guiado y Receptor Óptico Pasivo}
% =============================================================================

\begin{itemize}
    \item \textbf{Medio Óptico (Req. I.9):} Tramo de fibra óptica plástica (POF) cortado transversalmente a \SI{90}{\degree} con longitud exacta de:
    \begin{equation}
    L_{\text{fibra}} = \SI{45.0}{\centi\meter} \in [\SI{30.0}{\centi\meter}, \SI{60.0}{\centi\meter}]
    \end{equation}
    ubicado exactamente en el centro de la tolerancia normativa.
    \item \textbf{Receptor Óptico Pasivo (Req. J.10):} Fotorresistencia LDR GL5528 constituida de material semiconductor pasivo (sin polarización activa ni ganancia electrónica interna). La resistencia superficial responde a la ley de potencias fenomenológica \cite{rose1963}:
    \begin{equation}
    R_{\text{LDR}}(E_v) = R_0 \cdot E_v^{-\gamma}
    \label{eq:ldr_fotometria}
    \end{equation}
    En la simulación circuital y el modelo físico calibrado:
    \begin{equation}
    R_{\text{LDR}}(I_{\text{LED}}) = 1500 + 38\,500 \cdot \exp\left( -\frac{I_{\text{LED}}}{\SI{3}{\milli\ampere}} \right)\ \Omega
    \label{eq:ldr_modelo_spice}
    \end{equation}
    variando desde $R_{\text{LDR(osc)}} \approx \SI{40.0}{\kilo\ohm}$ en corte luminoso hasta $R_{\text{LDR(max)}} \approx \SI{1.577}{\kilo\ohm} \approx \SI{1.58}{\kilo\ohm}$ bajo la irradiación máxima del código $D=15$.
\end{itemize}

% =============================================================================
\section{Etapa 6: Generador de Onda Cuadrada (NE555 Astable)}
% =============================================================================

\subsection{Deducción Matemática a partir de la Ecuación Diferencial Transitoria RC}
El oscilador astable basado en el circuito integrado comercial NE555 opera mediante la carga y descarga cíclica del condensador $C_1 = \SI{100}{\nano\farad}$ a través de la resistencia fija $R_1 = \SI{1.0}{\kilo\ohm}$ y la fotorresistencia $R_{\text{LDR}}$ \cite{camenzind2005,ti_ne555}.

La solución analítica general de la ecuación diferencial de primer orden para la tensión en un condensador es:
\begin{equation}
v_C(t) = V_{\text{final}} + \left( V_{\text{inicial}} - V_{\text{final}} \right) e^{-t / \tau}
\label{eq:rc_transitorio}
\end{equation}

\paragraph{1. Semiciclo de Carga --- Tiempo en Alto ($t_H$):}
Durante este intervalo, el transistor interno de descarga (Pin 7) permanece en corte. El condensador se carga desde el umbral inferior $V_{\text{inicial}} = \frac{1}{3}V_{\text{CC}}$ hacia $V_{\text{final}} = V_{\text{CC}}$ a través de la resistencia equivalente de carga $(R_1 + R_{\text{LDR}})$, con constante de tiempo $\tau_c = (R_1 + R_{\text{LDR}}) C_1$. 

La conmutación ocurre cuando $v_C(t_H) = \frac{2}{3}V_{\text{CC}}$:
\begin{equation}
\frac{2}{3}V_{\text{CC}} = V_{\text{CC}} + \left( \frac{1}{3}V_{\text{CC}} - V_{\text{CC}} \right) e^{-t_H / \tau_c}
\end{equation}
\begin{equation}
-\frac{1}{3}V_{\text{CC}} = -\frac{2}{3}V_{\text{CC}} \, e^{-t_H / \tau_c} \implies e^{-t_H / \tau_c} = \frac{1}{2}
\end{equation}
Aplicando el logaritmo natural:
\begin{equation}
-\frac{t_H}{\tau_c} = \ln\left( \frac{1}{2} \right) = -\ln(2) \implies \boxed{t_H = \ln(2) \cdot (R_1 + R_{\text{LDR}}) \cdot C_1 \approx 0.69315\,(R_1 + R_{\text{LDR}})\,C_1}
\label{eq:th_deducido}
\end{equation}

\paragraph{2. Semiciclo de Descarga --- Tiempo en Bajo ($t_L$):}
Al alcanzar $\frac{2}{3}V_{\text{CC}}$, el flip-flop interno satura el transistor de descarga del Pin 7 conectándolo a tierra. El condensador se descarga desde $V_{\text{inicial}} = \frac{2}{3}V_{\text{CC}}$ hacia $V_{\text{final}} = \SI{0}{\volt}$ exclusivamente a través de $R_{\text{LDR}}$, con constante de tiempo $\tau_d = R_{\text{LDR}} C_1$.

La conmutación se activa cuando $v_C(t_L) = \frac{1}{3}V_{\text{CC}}$:
\begin{equation}
\frac{1}{3}V_{\text{CC}} = 0 + \left( \frac{2}{3}V_{\text{CC}} - 0 \right) e^{-t_L / \tau_d} \implies e^{-t_L / \tau_d} = \frac{1}{2}
\end{equation}
\begin{equation}
\boxed{t_L = \ln(2) \cdot R_{\text{LDR}} \cdot C_1 \approx 0.69315\,R_{\text{LDR}}\,C_1}
\label{eq:tl_deducido}
\end{equation}

\paragraph{3. Periodo y Frecuencia Fundamental de Oscilación ($f_{555}$):}
El periodo total de la forma de onda es:
\begin{equation}
T = t_H + t_L = \ln(2)(R_1 + R_{\text{LDR}})C_1 + \ln(2)R_{\text{LDR}}C_1 = \ln(2)\,(R_1 + 2R_{\text{LDR}})\,C_1
\label{eq:periodo_555}
\end{equation}
La frecuencia resultante es:
\begin{equation}
\boxed{f_{555} = \frac{1}{T} = \frac{1}{\ln(2)\,(R_1 + 2R_{\text{LDR}})\,C_1} = \frac{1.442695}{(R_1 + 2R_{\text{LDR}})\,C_1} \approx \frac{1.443}{(R_1 + 2R_{\text{LDR}})\,C_1}}
\label{eq:frecuencia_555}
\end{equation}

\paragraph{4. Ciclo de Trabajo (Duty Cycle):}
\begin{equation}
\mathrm{Duty\ Cycle} = \frac{t_H}{t_H + t_L} \times \SI{100}{\percent} = \frac{R_1 + R_{\text{LDR}}}{R_1 + 2R_{\text{LDR}}} \times \SI{100}{\percent}
\label{eq:duty_cycle}
\end{equation}

\subsection{Evaluación Numérica del Rango Dinámico Acústico}
Para $R_1 = \SI{1.0}{\kilo\ohm}$ y $C_1 = \SI{100}{\nano\farad}$:
\begin{enumerate}
    \item \textbf{Penumbra o Estado Apagado ($D \le 5$, $R_{\text{LDR}} = \SI{40.0}{\kilo\ohm}$):}
    \begin{equation}
    f_{\text{teo}} = \frac{1.443}{(1000 + 2 \times 40\,000)\times 10^{-7}} = \frac{1.443}{8.10 \times 10^{-3}} = \mathbf{\SI{178.1}{\hertz}}
    \end{equation}
    Simulación en LTspice: $f_{\text{sim}} = \SI{176.7}{\hertz}$ ($E_{\text{rel}} = \SI{0.81}{\percent}$, tono grave). Ciclo de trabajo: \SI{50.3}{\percent}.
    \item \textbf{Máxima Iluminación Óptica ($D = 15$, $R_{\text{LDR}} = \SI{1576.6}{\ohm} \approx \SI{1.58}{\kilo\ohm}$):}
    \begin{equation}
    f_{\text{teo}} = \frac{1.443}{(1000 + 2 \times 1576.6)\times 10^{-7}} = \frac{1.443}{4.1532 \times 10^{-4}} = \mathbf{\SI{3474.3}{\hertz}}
    \end{equation}
    *(Nota: Si se redondea $R_{\text{LDR}} = \SI{1.58}{\kilo\ohm}$, $f = 1.443 / (4160 \times 10^{-7}) = \SI{3468.8}{\hertz}$)*.
    Simulación en LTspice: $f_{\text{sim}} = \SI{3443.0}{\hertz}$ ($E_{\text{rel}} = \SI{0.90}{\percent}$, tono agudo). Ciclo de trabajo: \SI{62.0}{\percent}.
\end{enumerate}
Todo el rango de frecuencias reside plenamente en el espectro audible humano (\SI{20}{\hertz}--\SI{20}{\kilo\hertz}) y coincide con la máxima sensibilidad acústica del transductor piezoeléctrico.

% =============================================================================
\section{Etapa 7: Transductor Acústico Pasivo y Acoplamiento AC}
% =============================================================================

\subsection{Justificación del Desacoplo de Continua (Req. L.12)}
Un buzzer piezoeléctrico pasivo está constituido por una lámina cerámica de titanato-circonato de plomo (PZT) adherida a un disco elástico metálico \cite{cady1964}. La deformación mecánica periódica que genera presión sonora audible se induce exclusivamente por gradientes dinámicos de potencial eléctrico alterno ($\Delta V$).

La salida del temporizador 555 (Pin 3) conmuta entre $\approx \SI{0}{\volt}$ y $V_{\text{CC}}$, portando una componente de corriente continua no nula:
\begin{equation}
V_{\text{DC}} = \frac{1}{T}\int_0^T v_{\text{out}}(t)\,dt \approx \mathrm{Duty} \times V_{\text{CC}} \approx 0.50 \times \SI{5.0}{\volt} = \SI{2.50}{\volt}
\end{equation}
La presencia sostenida de una polarización continua genera fatiga dieléctrica, flexión asimétrica y calentamiento óhmico en el transductor cerámico pasivo \cite{cady1964}.

\subsection{Cálculo del Filtro Pasa-Altos de Acoplo AC}
Para eliminar completamente $V_{\text{DC}}$, se interpone en serie un condensador electrolítico radial de $C_{\text{acoplo}} = \SI{10.0}{\micro\farad}$. Modelando la impedancia del buzzer pasivo como una resistencia equivalente $R_L \approx \SI{1.0}{\kilo\ohm}$, el conjunto conforma un filtro pasa-altos pasivo RC de primer orden \cite{alexander2018}:
\begin{equation}
f_c = \frac{1}{2\pi \cdot R_L \cdot C_{\text{acoplo}}} = \frac{1}{2\pi \times (\SI{1000}{\ohm}) \times (\SI{10.0e-6}{\farad})} = \frac{1}{6.283185 \times 10^{-2}} = \mathbf{\SI{15.915}{\hertz}} \approx \mathbf{\SI{15.9}{\hertz}}
\label{eq:fc_acoplo}
\end{equation}

Dado que:
\begin{equation}
f_c = \SI{15.9}{\hertz} \ll f_{\text{min}} = \SI{176.7}{\hertz} \quad \left(\frac{f_{\text{min}}}{f_c} \approx 11.1\right)
\end{equation}
la frecuencia de corte se ubica más de una década por debajo del tono más grave del sistema, garantizando ganancia unitaria sin atenuación de amplitud en ninguna frecuencia útil y bloqueando al \SI{100}{\percent} el nivel continuo.

% =============================================================================
\section{Esquema de Alimentación y Cumplimiento Normativo (P.17 y Q.18)}
% =============================================================================

\subsection{Requerimientos Mandatorios de Alimentación}
Los requerimientos normativos \textbf{P.17} y \textbf{Q.18} establecen directrices infranqueables:
\begin{enumerate}
    \item \textbf{Desconexión Absoluta del Riel de \SI{+5}{\volt} del Arduino (Req. P.17):} El pin de \SI{+5}{\volt} del microcontrolador Arduino UNO no suministra alimentación a ninguna etapa analógica ni a los circuitos integrados de la protoboard.
    \item \textbf{Prohibición Expresa de Baterías Portátiles (Req. Q.18):} Queda terminantemente prohibido el empleo de baterías (AA, AAA o de 9V) durante la toma de datos cuantitativos finales y para la defensa evaluativa formal en semana 19. Su uso se restringió exclusivamente a una validación preliminar cualitativa de continuidad de banco tras concluir el ensamble físico.
    \item \textbf{Fuente de Alimentación DC Regulada de Laboratorio de Doble Canal:}
    \begin{itemize}
        \item \textbf{Canal 1 (Ajustado a \SI{+5.00}{\volt}):} Alimenta el temporizador NE555 (Pines 8 y 4).
        \item \textbf{Canal 2 (Ajustado a \SI{+9.00}{\volt}):} Alimenta el amplificador operacional LM358N (Pin 8).
        \item \textbf{Referencia Unificada (GND Común):} Los bornes negativos de ambos canales de la fuente de alimentación y el terminal \texttt{GND} del Arduino se unen en un único bus de tierra común (\SI{0}{\volt}).
    \end{itemize}
\end{enumerate}

\subsection{Diagnóstico Experimental de Laboratorio y Corrección de Puentes}
Durante la sesión de pruebas en el laboratorio de la Escuela de Ingeniería Electrónica del ITCR (efectuada por Mario Cerdas e Ian Fuentes), se registraron anomalías en las lecturas de tensión de salida de la escalera DAC. Tras auditar la protoboard y reproducir el comportamiento en TinkerCAD Circuits, se constató que los terminales de realimentación del encapsulado DIP-8 del LM358N se encontraban abiertos.

Se normalizó el circuito aplicando tres puentes rígidos directos:
\begin{enumerate}
    \item \textbf{Puente 1 (Pin 1 a Pin 2):} Cierra el lazo de realimentación negativa del OpAmp 1A (Buffer seguidor unitario del DAC R--2R, Etapa 2).
    \item \textbf{Puente 2 (Pin 6 a Pin 7):} Cierra el lazo de realimentación negativa del OpAmp 1B (Buffer seguidor unitario de la interfaz, Etapa 3).
    \item \textbf{Enlace Inter-etapas en U (Pin 1 a Pin 5):} Transfiere la tensión analógica $V_{\text{DAC}}$ hacia la entrada no inversora de la Etapa 3 sin efecto de carga.
\end{enumerate}
Las pruebas de banco posteriores ratificaron la total linealidad monótona de la escalera de 16 escalones (\SI{0}{\volt} a \SI{4.6875}{\volt}).

% =============================================================================
\section{Cálculos Complementarios de Ingeniería (Para Discusión del Equipo)}
% =============================================================================

A continuación se presentan cuatro análisis cuantitativos avanzados desarrollados para enriquecer la memoria técnica y proporcionar argumentos rigurosos durante la defensa oral ante el profesor. El equipo de trabajo puede evaluar cuáles incorporar íntegramente en la sustentación y cuáles mantener como respaldo de diseño:

\subsection{Análisis de Tolerancias en Peor Caso y No Linealidad (DNL e INL)}
El requerimiento \textbf{O.16} estipula un error relativo inferior al \SI{10}{\percent}. Empleando resistores comerciales de película metálica con tolerancia del \SI{1}{\percent} ($\epsilon = \pm 0.01$), se evalúa la no linealidad diferencial (DNL) en la transición más crítica de la escalera R--2R: el cruce de media escala ($0111_2 \to 1000_2$):
\begin{itemize}
    \item En $D=7$ ($0111_2$), la salida depende de la suma de los bits 0, 1 y 2 con peso conjunto $7/16$.
    \item En $D=8$ ($1000_2$), se desconectan los tres bits inferiores y conmuta exclusivamente el bit MSB ($b_3$) con peso nominal $1/2 = 8/16$.
\end{itemize}

Bajo la condición de dispersión en peor caso (rama del bit 3 desviada en \SI{-1}{\percent} y ramas inferiores en \SI{+1}{\percent}):
\begin{equation}
\Delta V_{\text{paso(peor)}} = V_{\text{LSB}} \pm 2\epsilon \cdot V_{\text{ref}} = \SI{0.3125}{\volt} \pm 2(0.01)(\SI{5.0}{\volt}) = \SI{0.3125}{\volt} \pm \SI{0.1000}{\volt}
\end{equation}
Expresando el error de salto en unidades de LSB:
\begin{equation}
\mathrm{DNL}_{\text{max}} = \frac{\SI{0.1000}{\volt}}{\SI{0.3125}{\volt}} = \mathbf{0.32\text{ LSB}} < \mathbf{1.0\text{ LSB}}
\label{eq:dnl_calculo}
\end{equation}
\textbf{Conclusión de Ingeniería:} Dado que $|\mathrm{DNL}| < 1.0\text{ LSB}$, \textbf{la monotonicidad de la escalera está matemáticamente garantizada} (no existen códigos perdidos ni inversiones de pendiente). El error integral máximo (INL) es:
\begin{equation}
\mathrm{INL}_{\text{max}} \approx \epsilon \cdot \frac{V_{\text{ref}}}{2} = 0.01 \times \SI{2.50}{\volt} = \SI{25.0}{\milli\volt} \implies E_{\text{rel(peor)}} \approx \frac{\SI{25}{\milli\volt}}{\SI{2500}{\milli\volt}} \times \SI{100}{\percent} = \mathbf{\SI{1.0}{\percent}} \ll \mathbf{\SI{10}{\percent}}
\end{equation}

\subsection{Balance de Potencia Total del Sistema y Elevación Térmica}
Para dimensionar adecuadamente la Fuente DC Regulada de Laboratorio y verificar la estabilidad térmica de los componentes activos:
\begin{enumerate}
    \item \textbf{Canal 1 (\SI{+5.00}{\volt} --- NE555 y temporización):}
    \begin{itemize}
        \item Corriente de reposo del NE555: $I_{\text{CC(555)}} \approx \SI{6.0}{\milli\ampere}$ \cite{ti_ne555}.
        \item Corriente promedio por $R_1$: $I_{R1} \approx \frac{V_{\text{CC}} - \frac{2}{3}V_{\text{CC}}}{R_1} \times \mathrm{Duty} \approx \frac{\SI{1.67}{\volt}}{\SI{1000}{\ohm}} \times 0.60 \approx \SI{1.0}{\milli\ampere}$.
        \item Consumo total en CH1: $I_{\text{CH1}} \approx \SI{7.0}{\milli\ampere} \implies P_{\text{CH1}} = \SI{5.0}{\volt} \times \SI{7.0}{\milli\ampere} = \mathbf{\SI{35.0}{\milli\watt}}$.
    \end{itemize}
    \item \textbf{Canal 2 (\SI{+9.00}{\volt} --- LM358N y Driver LED):}
    \begin{itemize}
        \item Corriente de reposo del LM358N (ambos operacionales): $I_{\text{CC(LM358)}} \approx \SI{1.2}{\milli\ampere}$ \cite{ti_lm358}.
        \item Corriente máxima entregada al LED en $D=15$: $I_{\text{LED(max)}} = \SI{18.66}{\milli\ampere}$.
        \item Consumo total máximo en CH2: $I_{\text{CH2(max)}} = 1.2 + 18.66 = \SI{19.86}{\milli\ampere} \implies P_{\text{CH2(max)}} = \SI{9.0}{\volt} \times \SI{19.86}{\milli\ampere} = \mathbf{\SI{178.7}{\milli\watt}}$.
    \end{itemize}
    \item \textbf{Potencia total suministrada:} $P_{\text{total}} = \SI{35.0}{\milli\watt} + \SI{178.7}{\milli\watt} = \mathbf{\SI{213.7}{\milli\watt}} \approx \mathbf{\SI{0.21}{\watt}}$.
    \item \textbf{Disipación interna y elevación de temperatura en el LM358N (DIP-8):}
    En $D=15$, la caída de tensión en el transistor de salida del OpAmp 1B es $V_{\text{CE}} = V_{\text{OP}} - V_{\text{interfaz}} = \SI{9.00}{\volt} - \SI{4.6875}{\volt} = \SI{4.3125}{\volt}$.
    \begin{equation}
    P_{\text{chip}} = (\SI{4.3125}{\volt} \times \SI{18.66}{\milli\ampere}) + (\SI{9.0}{\volt} \times \SI{1.2}{\milli\ampere}) = \SI{80.5}{\milli\watt} + \SI{10.8}{\milli\watt} = \mathbf{\SI{91.3}{\milli\watt}}
    \end{equation}
    Con una resistencia térmica unión-ambiente típica de encapsulado DIP-8 $\theta_{\text{JA}} \approx \SI{100}{\celsius\per\watt}$ \cite{ti_lm358}:
    \begin{equation}
    \Delta T = P_{\text{chip}} \cdot \theta_{\text{JA}} = \SI{0.0913}{\watt} \times \SI{100}{\celsius\per\watt} = \mathbf{\SI{9.1}{\celsius}} \implies T_J \approx \SI{25}{\celsius} + \SI{9.1}{\celsius} = \mathbf{\SI{34.1}{\celsius}} \ll \SI{150}{\celsius}
    \end{equation}
    El amplificador opera totalmente frío, garantizando vida útil prolongada sin necesidad de disipador.
\end{enumerate}

\subsection{Balance de Enlace Óptico (Optical Link Budget) y Pérdidas de Fresnel}
Para caracterizar cuantitativamente la atenuación del canal óptico guiado de \SI{45}{\centi\meter}:
\begin{enumerate}
    \item \textbf{Pérdidas por Reflexión de Fresnel en las Facetas de la Fibra:}
    En la transición aire--núcleo de PMMA ($n_{\text{aire}} = 1.000$, $n_1 = 1.492$) con incidencia normal \cite{tipler2010,ziemann2008}:
    \begin{equation}
    R = \left( \frac{n_1 - n_{\text{aire}}}{n_1 + n_{\text{aire}}} \right)^2 = \left( \frac{1.492 - 1.000}{1.492 + 1.000} \right)^2 = \left( \frac{0.492}{2.492} \right)^2 = 0.03898 \approx \mathbf{\SI{3.90}{\percent}}
    \end{equation}
    Pérdida en decibeles por faceta: $A_{\text{Fresnel}} = -10\log_{10}(1 - R) = -10\log_{10}(0.9610) = \mathbf{\SI{0.173}{\decibel}/\text{faceta}}$. Para las dos caras (entrada y salida):
    \begin{equation}
    A_{\text{Fresnel(total)}} = 2 \times \SI{0.173}{\decibel} = \mathbf{\SI{0.346}{\decibel}}
    \end{equation}
    \item \textbf{Eficiencia de Acoplamiento Geométrico (Cono de Aceptación):}
    Asumiendo emisión lambertiana del LED, la fracción de potencia capturada por la fibra está acotada por su apertura numérica \cite{kasap2013}:
    \begin{equation}
    \eta_{\text{acoplo}} \approx \mathrm{NA}^2 = (0.510)^2 = 0.2601 \approx \mathbf{\SI{26.0}{\percent}} \implies A_{\text{acoplo}} = -10\log_{10}(0.2601) = \mathbf{\SI{5.85}{\decibel}}
    \end{equation}
    \item \textbf{Atenuación Distribuida en la Fibra POF ($L = \SI{0.45}{\meter}$):}
    \begin{equation}
    A_{\text{fibra}} = \alpha \cdot L_{\text{fibra}} = \SI{0.15}{\decibel\per\meter} \times \SI{0.45}{\meter} = \mathbf{\SI{0.0675}{\decibel}}
    \end{equation}
    \item \textbf{Atenuación Óptica Total del Enlace:}
    \begin{equation}
    A_{\text{total}} = A_{\text{acoplo}} + A_{\text{Fresnel(total)}} + A_{\text{fibra}} = \SI{5.85}{\decibel} + \SI{0.35}{\decibel} + \SI{0.07}{\decibel} = \mathbf{\SI{6.27}{\decibel}}
    \end{equation}
    Para una potencia óptica emitida por el LED de $P_{\text{tx}} \approx \SI{2.50}{\milli\watt}$ (\SI{+3.98}{\decibel m}), la potencia incidente en la LDR es:
    \begin{equation}
    P_{\text{rx}} = \SI{+3.98}{\decibel m} - \SI{6.27}{\decibel} = \SI{-2.29}{\decibel m} \approx \mathbf{\SI{0.59}{\milli\watt}} = \SI{590}{\micro\watt}
    \end{equation}
    Esta irradiancia es óptima para reducir la resistencia del sulfuro de cadmio hasta \SI{1.58}{\kilo\ohm}.
\end{enumerate}

\subsection{Análisis Dinámico: Slew Rate del OpAmp vs. Cinética Fotoconductora de la LDR}
Para justificar técnicamente el tiempo de permanencia mínimo ($T_{\text{est}} \ge \SI{1}{\milli\second}$) exigido en el requerimiento \textbf{C.3}:
\begin{enumerate}
    \item \textbf{Tiempo de Respuesta del DAC y OpAmp LM358N:}
    La constante de tiempo parásita de la escalera R--2R con la capacitancia de entrada del operacional ($C_{\text{in}} \approx \SI{3}{\pico\farad}$) es $\tau_{\text{DAC}} = \SI{10}{\kilo\ohm} \times \SI{3}{\pico\farad} = \SI{30}{\nano\second}$. La velocidad de respuesta del LM358N está gobernada por su Slew Rate ($\mathrm{SR} = \SI{0.50}{\volt\per\micro\second}$) \cite{ti_lm358}. El tiempo de transición para una excursión completa ($\SI{0}{\volt} \to \SI{4.6875}{\volt}$) es:
    \begin{equation}
    t_{\text{SR}} = \frac{\Delta V}{\mathrm{SR}} = \frac{\SI{4.6875}{\volt}}{\SI{0.50}{\volt\per\micro\second}} = \mathbf{\SI{9.375}{\micro\second}}
    \end{equation}
    En comparación con el tiempo de retención mínimo permitido ($T_{\text{est}} = \SI{1}{\milli\second} = \SI{1000}{\micro\second}$):
    \begin{equation}
    \frac{t_{\text{SR}}}{T_{\text{est(min)}}} = \frac{\SI{9.375}{\micro\second}}{\SI{1000}{\micro\second}} \times \SI{100}{\percent} = \mathbf{\SI{0.94}{\percent}}
    \end{equation}
    La tensión analógica se establece en menos del \SI{1}{\percent} del intervalo temporal de la muestra.
    \item \textbf{Cinética Fotoconductora del Sulfuro de Cadmio (CdS):}
    A diferencia de la respuesta electrónica del DAC ($< \SI{10}{\micro\second}$), el proceso de recombinación de portadores libres en la fotorresistencia GL5528 está limitado por estados trampa profundos \cite{rose1963}, exhibiendo un tiempo de subida $t_r \approx \SI{20}{\milli\second}$--\SI{30}{\milli\second} y un tiempo de caída (retorno a oscuridad) $t_f \approx \SI{30}{\milli\second}$--\SI{40}{\milli\second}.
    
    \textbf{Conclusión de Ingeniería:} A frecuencias de conmutación muy elevadas ($T_{\text{est}} \approx \SI{1}{\milli\second}$), la LDR actúa como un filtro pasa-bajos promediador. Para lograr que la frecuencia del tono acústico del 555 alcance plenamente el estado estacionario en cada muestra, el tiempo de permanencia óptimo de operación es $T_{\text{est}} \ge \SI{500}{\milli\second}$ (habiéndose configurado \SI{1500}{\milli\second} por defecto en el firmware de Arduino).
\end{enumerate}

% =============================================================================
\section{Cuadro Maestro de Resultados: Teórico vs. LTspice vs. Experimental}
% =============================================================================

El requerimiento \textbf{O.16} estipula que el error de medición del sistema conversor debe ser estrictamente inferior al \SI{10}{\percent}:
\begin{equation}
E_{\text{rel}}(\%) = \left| \frac{V_{\text{sim}} - V_{\text{teo}}}{V_{\text{teo}}} \right| \times \SI{100}{\percent} < \SI{10.0}{\percent}
\label{eq:error_rel_def}
\end{equation}

\begin{table}[htbp]
\centering
\caption{Consolidado Teórico, Simulado en LTspice y Plantilla de Laboratorio para los 16 Estados.}
\label{tab:maestro_resultados}
\small
\renewcommand{\arraystretch}{1.12}
\setlength{\tabcolsep}{4.5pt}
\begin{tabular}{ccccccccc}
\toprule
\textbf{Dec ($D$)} & \textbf{Binario} & $V_{\text{teo}}$ \textbf{(V)} & $V_{\text{sim}}$ \textbf{(V)} & $V_{\text{exp}}$ \textbf{(V)} & \textbf{Error Sim (\%)} & $I_{\text{LED}}$ \textbf{(mA)} & $R_{\text{LDR}}$ \textbf{(k$\Omega$)} & $f_{555}$ \textbf{(Hz)} \\
\midrule
0  & 0000 & 0.0000 & 0.0050 & -- & Sat. $V_{\text{OL}}$ & 0.00 & 40.00 & 176.7 \\
1  & 0001 & 0.3125 & 0.3125 & -- & \SI{0.0012}{\percent} & 0.00 & 40.00 & 176.7 \\
2  & 0010 & 0.6250 & 0.6250 & -- & \SI{0.0000}{\percent} & 0.00 & 40.00 & 176.7 \\
3  & 0011 & 0.9375 & 0.9375 & -- & \SI{0.0004}{\percent} & 0.00 & 40.00 & 176.7 \\
4  & 0100 & 1.2500 & 1.2500 & -- & \SI{0.0006}{\percent} & 0.00 & 40.00 & 176.7 \\
5  & 0101 & 1.5625 & 1.5625 & -- & \SI{0.0007}{\percent} & 0.03 & 39.57 & 178.6 \\
6  & 0110 & 1.8750 & 1.8750 & -- & \SI{0.0008}{\percent} & 1.04 & 28.73 & 244.9 \\
7  & 0111 & 2.1875 & 2.1875 & -- & \SI{0.0009}{\percent} & 2.81 & 16.68 & 416.6 \\
8  & 1000 & 2.5000 & 2.5000 & -- & \SI{0.0009}{\percent} & 4.69 &  9.56 & 710.9 \\
9  & 1001 & 2.8125 & 2.8125 & -- & \SI{0.0009}{\percent} & 6.63 &  5.71 & 1151.9 \\
10 & 1010 & 3.1250 & 3.1250 & -- & \SI{0.0010}{\percent} & 8.61 &  3.68 & 1711.6 \\
11 & 1011 & 3.4375 & 3.4375 & -- & \SI{0.0010}{\percent} & 10.60 & 2.62 & 2291.0 \\
12 & 1100 & 3.7500 & 3.7500 & -- & \SI{0.0010}{\percent} & 12.61 & 2.08 & 2776.7 \\
13 & 1101 & 4.0625 & 4.0624 & -- & \SI{0.0010}{\percent} & 14.62 & 1.79 & 3116.4 \\
14 & 1110 & 4.3750 & 4.3749 & -- & \SI{0.0010}{\percent} & 16.64 & 1.65 & 3325.0 \\
15 & 1111 & 4.6875 & 4.6874 & -- & \SI{0.0010}{\percent} & 18.66 & 1.58 & 3443.0 \\
\bottomrule
\end{tabular}
\end{table}

\noindent \textbf{Análisis de Validación del Cuadro~\ref{tab:maestro_resultados}:}
\begin{enumerate}
    \item Para todos los estados activos ($D=1$ a $15$), el error relativo simulado es menor al \SI{0.0012}{\percent}, superando en cuatro órdenes de magnitud el requerimiento del \SI{10}{\percent}.
    \item En el estado $D=0$ ($0000_2$), la salida teórica es \SI{0.0000}{\volt}. La simulación arroja un offset de \SI{5.0}{\milli\volt} correspondiente a la tensión residual de saturación inferior $V_{\text{OL}}$ del LM358 con fuente unipolar.
    \item La corriente del LED se mantiene estrictamente por debajo de los \SI{20}{\milli\ampere} ($I_{\text{max}} = \SI{18.66}{\milli\ampere}$ en $D=15$), asegurando la integridad del semiconductor.
\end{enumerate}

% =============================================================================
\section{Referencias Bibliográficas Académicas (Normadas IEEE)}
% =============================================================================

\begin{thebibliography}{14}

\bibitem{alexander2018}
C.~K.~Alexander and M.~N.~O.~Sadiku, \emph{Fundamentos de circuitos eléctricos}, 6ta~ed. México D.F.: McGraw-Hill Interamericana, 2018.

\bibitem{hayt2012}
W.~H.~Hayt, J.~E.~Kemmerly, and S.~M.~Durbin, \emph{Análisis de circuitos en ingeniería}, 8va~ed. México D.F.: McGraw-Hill, 2012.

\bibitem{sedra2020}
A.~S.~Sedra, K.~C.~Smith, T.~C.~Carusone, and V.~Gaudet, \emph{Microelectronic Circuits}, 8th~ed. New York: Oxford University Press, 2020.

\bibitem{franco2015}
S.~Franco, \emph{Design with Operational Amplifiers and Analog Integrated Circuits}, 4th~ed. New York: McGraw-Hill Education, 2015.

\bibitem{proakis2007}
J.~G.~Proakis and D.~G.~Manolakis, \emph{Digital Signal Processing: Principles, Algorithms, and Applications}, 4th~ed. Upper Saddle River, NJ: Pearson Prentice Hall, 2007.

\bibitem{kasap2013}
S.~O.~Kasap, \emph{Optoelectronics and Photonics: Principles and Practices}, 2nd~ed. Boston: Pearson, 2013.

\bibitem{tipler2010}
P.~A.~Tipler y G.~Mosca, \emph{Física para la ciencia y la tecnología. Vol. 2: Electricidad y magnetismo / Luz}, 6ta~ed. Barcelona: Editorial Reverté, 2010.

\bibitem{ziemann2008}
O.~Ziemann, J.~Krauser, P.~E.~Zamzow, and W.~Daum, \emph{POF Handbook: Optical Short Range Transmission Systems}, 2nd~ed. Berlin Heidelberg: Springer-Verlag, 2008.

\bibitem{rose1963}
A.~Rose, \emph{Concepts in Photoconductivity and Allied Problems}. New York: Interscience Publishers, John Wiley \& Sons, 1963.

\bibitem{streetman2016}
B.~G.~Streetman and S.~K.~Banerjee, \emph{Solid State Electronic Devices}, 7th~ed. Upper Saddle River, NJ: Pearson, 2016.

\bibitem{camenzind2005}
H.~R.~Camenzind, \emph{Designing Analog Chips}. Virtualbookworm Publishing, 2005.

\bibitem{cady1964}
W.~G.~Cady, \emph{Piezoelectricity: An Introduction to the Theory and Applications of Electromechanical Phenomena in Crystals}. New York: Dover Publications, 1964.

\bibitem{ti_lm358}
Texas Instruments, \emph{LMx58-N Low-Power, Dual-Operational Amplifiers Datasheet}, Doc.~SNOSBT3J, Dallas, TX, 2023.

\bibitem{ti_ne555}
Texas Instruments, \emph{NA555, NE555, SA555, SE555 Precision Timers Datasheet}, Doc.~SLFS022I, Dallas, TX, 2022.

\end{thebibliography}

\end{document}